\begin{proof}

\begin{enumerate}[label={$\langle1\rangle$\arabic*.},leftmargin=*]
\item Shrink $B$ so that the holomorphic Lie bundle of $G$ has a
frame $e_1,\ldots,e_n$.  Its period local system has a basis
of holomorphic sections
\[
 \lambda_1(b),\ldots,\lambda_{2n}(b)\in\C^n,
 \qquad
 G_b=\C^n/\sum_{a=1}^{2n}\Z\lambda_a(b).
\]
Indeed the lattice is locally constant over the simply
connected bidisc, and the local inverses of the relative
exponential map make its sections holomorphic.  The periods
are a real basis at every $b$.

\item The infinitesimal torsor action gives holomorphic vector
fields $V_1,\ldots,V_n$ on $X_U$.  Flatness makes
$Z=g^{-1}(0)$ have complex codimension two in $X$.
Hartogs extension in local coordinates therefore extends
each coefficient of each $V_j$ uniquely across $Z$; uniqueness
glues these local extensions to holomorphic vector fields
on $X$.  This use of Hartogs applies to functions on the
complement of an analytic set of codimension at least two
in a complex manifold, hence also coefficientwise to any
locally free sheaf.  It does not assert extension across a
divisor.  The identities
\[
 dg(V_j)=0,\qquad [V_i,V_j]=0
\]
hold on the dense open subset $X_U$ and consequently on $X$.

\item Every $V_j$, and every linear combination with coefficients
pulled back from $B$, is complete.  Its integral curve stays
in a compact fibre; a finite endpoint of a maximal solution
would have an accumulation point in that fibre, where the
holomorphic ordinary differential equation extends the
solution.  The flows thus define a holomorphic action of
$\C^n$ on $X$ over $B$.

\item For $a=(a_1,\ldots,a_{2n})\in\R^{2n}$ take the time-one
flow of the vertical holomorphic vector field whose Lie
coordinate is
\[
 \sum_{j=1}^{2n}a_j\lambda_j(g(x)).
\]
These maps depend smoothly on $a$ and add under composition:
all their coefficients are constant along every vertical
flow.  For $a\in\Z^{2n}$ the map is the identity on $X_U$,
hence everywhere by holomorphic uniqueness.  We obtain an
action of the fixed compact real torus
$T=\R^{2n}/\Z^{2n}$ by holomorphic automorphisms of $X$,
over $B$.  On the fibre over $b$, it is the holomorphic
action of $G_b$ obtained by factoring the $\C^n$ action by
its period lattice.  The action of $T$ is free on $X_U$.
The variable period identities also make the holomorphic
Lie-bundle action descend to a holomorphic action
$G\times_BX\to X$.  Its restriction therefore acts on
each scheme-theoretic fibre, including its nilpotent
structure.

\item We show that its stabilizers have dimension zero everywhere.
Let $x\in X$, $b=g(x)$, and let $H$ be the identity component
of the stabilizer.  Viewed in $G_b$, it is a compact complex
subgroup: it is the identity component of the inverse image
of $x$ under the holomorphic orbit map $G_b\to X_b$.
Put
\[
 K_x=\ker(dg_x)=T_x(X_b),
\]
where the right-hand side is the Zariski tangent space of
the scheme-theoretic fibre.  Differentiating the holomorphic
$G_b$ action on this fibre gives a holomorphic representation
$H\to\operatorname{GL}(K_x)$.  It is trivial, because every
holomorphic function on the compact connected complex group
$H$ is constant.  The induced representation on
$T_xX/K_x$ is also trivial: this quotient injects into
$T_bB$ by $dg_x$, and every group element acts over $B$.
Averaging a Hermitian inner product over the compact group
$H$ gives an invariant complement to $K_x$.  The full
derivative representation on $T_xX$ is therefore trivial.

\item Here the required local linearization has a direct proof.
Choose holomorphic coordinates $z$ centred at $x$ and a
sufficiently small $H$-invariant neighbourhood on which all
$z(hy)$ are defined.  Such a neighbourhood exists by
compactness of $H$ and continuity of the action.  With
normalized Haar measure set
\[
 \Psi(y)=\int_H z(hy)\,dh.
\]
This map is holomorphic, invariant under $H$, and has
derivative the identity at $x$.  It is a holomorphic
coordinate system after shrinking.  Its invariance then
forces every $h\in H$ to be the identity near $x$, and hence
on connected $X$.  This contradicts freeness on $X_U$ unless
$H$ is trivial.  Thus every stabilizer is finite, and the
$n$ vector fields $V_j$ are linearly independent everywhere.

\item Their commuting flows give local holomorphic coordinates
$(z,s)\in\C^n\times\C^2$ in which they span the $z$
directions: choose a transverse complex surface and apply
the holomorphic inverse function theorem to the flow map.
In these coordinates $g$ has the form
\[
 g(z,s)=h(s).
\]
Equidimensionality implies that the germ $h^{-1}(b)$ is
zero-dimensional.  Consequently the germ
$h:(\C^2,0)\to(B,b)$ is finite; this is the local finiteness
criterion obtained by Weierstrass division for an isolated
fibre.  If $dg_x$ failed to have rank two, the determinant
of $dh$ would vanish at $s=0$.  It is not identically zero,
since the map is a submersion away from $Z$.  Its zero set
would therefore contain a complex curve through $0$.
Finiteness sends that curve to a curve of critical values.
This contradicts smoothness over $B\setminus\{0\}$.
Hence $g$ is a submersion everywhere.

\item The $G_b$ orbits now have the full dimension of every smooth
connected fibre and are compact.  Each fibre is a single
orbit.  Its finite stabilizer acts trivially on its tangent
space, because $G_b$ is abelian; it also acts trivially on
the base tangent quotient.  Averaging as above, now over
that finite group, again makes its full derivative trivial
and then makes the action locally trivial.  Freeness on
$X_U$ forces the stabilizer to be trivial.  Thus the extended
action is free and transitive on every fibre.  Local
holomorphic sections of the submersion identify $X$ with
$G$ over small base opens, which proves the torsor assertion.
The extension is unique because it agrees with the original
one on the dense open subset $X_U$.

\end{enumerate}
\end{proof}
